% ─────────────────────────────────────────────────
%  CHAPITRE 1 – PRESENTATIONS GENERALES
% ─────────────────────────────────────────────────
\chapter{Présentations générales}
\markboth{Présentations générales}{Présentations générales}

\section*{Introduction}
\addcontentsline{toc}{section}{Introduction}

Ce chapitre situe le projet dans son cadre général. Il présente d'abord
Innov4Africa, structure d'accueil du stage, afin de préciser l'environnement
professionnel dans lequel le travail a été réalisé. Il expose ensuite le
contexte du sujet, la problématique retenue, les objectifs visés, les
caractéristiques du projet et la démarche de travail adoptée.

% ─────────────────────────────────────────────────
\section{Présentation de la structure d'accueil}

\subsection{Présentation de Innov4Africa}
INNOV4AFRICA, une société anonyme au capital de dix (10) millions de FCFA, est située à la villa
numéro 9405, Sacré-Cœur 3, Ancienne Piste X VDN à Dakar.
L'objectif principal de INNOV4AFRICA est de contribuer au développement économique et social
de l'Afrique en stimulant l'innovation à travers les nouvelles technologies, couvrant divers secteurs
économiques. En tant qu'acteur de l'innovation, l'entreprise œuvre pour l'intégration des technologies
modernes afin de transformer et dynamiser l'ensemble du continent africain.
L'organisation d'INNOV4AFRICA repose sur une structure de holdings regroupant des filiales
spécialisées et des représentants commerciaux présents dans de nombreux pays africains, tant
francophones qu'anglophones, afin de couvrir l'ensemble du marché africain.
Le Président Directeur Général (DG) de INNOV4AFRICA se prénomme M. Khassoum WONE. Ce
dernier est entouré par un conseil d'administration d'une expertise reconnue regroupant des
personnalités indépendantes. M. Khassoum WONE est un ingénieur de conception en Informatique.
Il a été DG de l'Agence de l'Informatique De l'Etat (ADIE) du 27 avril 2012 au 23 juillet 2014. Il a
managé une équipe d'une centaine d'ingénieurs informaticiens/télécoms et a pris le pilotage de projets
d'envergure à gros budget au Sénégal.

\subsection{Domaines d'activités}

INNOV4AFRICA se distingue par sa capacité à concevoir et à déployer des innovations majeures dans le domaine des technologies de l'information et de la communication (TIC), au service des gouvernements, des entreprises et des citoyens à travers l'Afrique.

Le numérique d'INNOV4AFRICA touche plusieurs secteurs. L'entreprise est présente dans les services publics --- e-Government, e-Santé, e-Éducation --- et dans le soutien à l'innovation, où elle incube des start-ups et gère des fonds d'investissement. Côté technologie, son terrain va de la cybersécurité à la biométrie, en passant par les paiements électroniques et le transfert de technologies. Elle conçoit aussi des logiciels et des applications mobiles, et intervient sur les télécommunications et le e-commerce. Le reste de son activité tourne autour de la communication et de l'accompagnement : médias, systèmes d'identification, développement d'affaires, partenariats stratégiques.
% ─────────────────────────────────────────────────
\section{Présentation du sujet}

\subsection{Contexte}

Le secteur bancaire connaît une transformation profonde portée par la
digitalisation des services, l'ouverture des systèmes par API, l'exigence
d'instantanéité et l'arrivée de l'intelligence artificielle dans les processus
métiers. Les clients attendent désormais une expérience simple, continue et
personnalisée, tandis que les établissements financiers doivent maîtriser leurs
risques, respecter la conformité et optimiser leurs coûts d'exploitation.

Dans ce contexte, un Système d'Information Bancaire (SIB) ne sert plus seulement
à tenir les comptes. Il doit aujourd'hui réunir au même endroit la relation
client, les opérations bancaires et l'investissement, sans baisser la garde sur
la sécurité et la traçabilité. C'est pourquoi les banques, comme les autres
établissements financiers, ont besoin de solutions plus souples, plus ouvertes,
et capables d'intégrer de l'intelligence artificielle.

\textbf{I-SIB}, porté par INNOV4AFRICA, s'inscrit dans cette évolution. Il
vise à proposer un système d'information bancaire complet, construit autour
d'une architecture microservices, d'une logique full API, d'une disponibilité
24/7 et d'un usage de l'intelligence artificielle dans les modules métiers.

Le périmètre traité concerne deux modules complémentaires : le module
\textbf{RELATION CLIENT}, qui centralise la connaissance client et les
interactions commerciales, et le module \textbf{SGI}, qui permet d'intégrer la
gestion d'investissement directement dans l'application bancaire.

Pour faciliter la compréhension de ce livrable, les principaux termes et
concepts utilisés sont définis ci-dessous.

\begin{description}
    \item[I-SIB] Solution de SIB intelligente proposée par INNOV4AFRICA,
    intégrant les processus bancaires, l'IA, les API et une architecture
    microservices.
    \item[Microservice] Application autonome dédiée à une responsabilité métier
    précise, pouvant évoluer et être déployée indépendamment.
    \item[Portefeuille] Les actifs financiers que détient un client, par
    exemple des actions et des obligations.
    \item[Cotation] Le prix auquel se négocie un titre (action ou obligation)
    à un moment donné.
    \item[Passage d'ordres] L'opération par laquelle un client demande
    d'acheter ou de vendre un titre.
\end{description}

% ─────────────────────────────────────────────────
\subsection{Problématique}

Dans beaucoup d'environnements bancaires, la relation client et les services
d'investissement sont traités dans des outils séparés. Le client consulte ses
comptes dans une application bancaire, puis doit souvent utiliser un autre canal
ou contacter une autre structure pour accéder aux services de placement, acheter
ou vendre des actions BRVM \cite{brvm_actions}, suivre ses obligations ou visualiser la performance
de son portefeuille-titres.

Cette séparation limite la prise de décision du client, car même lorsqu'il peut
alimenter son compte d'investissement depuis son compte bancaire, il ne dispose
pas toujours d'une vue unifiée reliant sa liquidité bancaire, son
portefeuille-titres, les opportunités de marché et son historique relationnel
avec la banque. L'expérience est fragmentée, ce qui freine l'engagement sur les
produits d'investissement.

Pour les équipes internes, l'absence de passerelle entre les deux systèmes
empêche de croiser les données bancaires et les données d'investissement d'un
même client. Les interactions sont journalisées séparément dans des systèmes
sans communication entre eux, rendant l'historique incomplet. La banque n'a
aucune visibilité sur l'activité d'investissement du client, et la SGI ne
dispose d'aucun contexte bancaire. Le client sort complètement du périmètre de
l'établissement pour investir, en utilisant ses données bancaires dans un cadre
totalement distinct.

% ─────────────────────────────────────────────────
\subsection{Objectifs}

\subsubsection{Objectif général}

Disposer d'un système d'information bancaire intelligent et centralisé,
intégrant les modules \textbf{RELATION CLIENT} et \textbf{SGI}, permettant
d'améliorer l'expérience client, le suivi des portefeuilles et l'accès aux
produits d'investissement depuis une plateforme unique.

\subsubsection{Objectifs spécifiques}

Concrètement, ces objectifs couvrent plusieurs axes. Sur le plan du pilotage, le
projet vise à mettre en place un tableau de bord Relation Client adapté aux
besoins du client et du chargé de clientèle. Il doit aussi faciliter le suivi des
prospects dans un pipeline de conversion, avec un score attribué automatiquement
pour aider à identifier les profils les plus prioritaires. Au-delà, le système
doit offrir une vue 360° du client --- identité, KYC, comptes, produits, transactions
et interactions ---, dématérialiser les documents et courriers administratifs
liés à son cycle de vie, et intégrer les exigences KYC directement dans le
suivi. L'intelligence artificielle vient en appui pour repérer les risques de
départ et renforcer le suivi relationnel. Enfin, du côté investissement, la
plateforme donne accès aux portefeuilles (performance, gains et pertes,
répartition des actifs) et prend en charge les cotations, le passage d'ordres
ainsi que le suivi des achats et ventes de titres.

% ─────────────────────────────────────────────────
% ─────────────────────────────────────────────────
\section{Démarche méthodologique}

La démarche retenue est itérative et incrémentale. Ce choix est cohérent avec
la nature du projet. Dans cette section, nous allons présenter les différentes
caractéristiques du projet .

% ─────────────────────────────────────────────────

\subsection{Caractéristiques du projet}

Dans le cadre de ce stage, le travail porte sur le projet \textbf{I-SIB – Système d'Information Bancaire intelligent}, porté par l'entreprise \textbf{INNOV4AFRICA}. Il s'inscrit dans un contexte de transformation numérique des systèmes bancaires, avec un accent particulier sur les modules \textbf{RELATION CLIENT} et \textbf{SGI (Société de Gestion et d'Intermédiation)}.

Les principaux bénéficiaires de cette solution sont les banques, les institutions de microfinance, les établissements de paiement, les SGI, ainsi que les clients finaux et les chargés de clientèle. Le produit développé est une application bancaire modulaire, intelligente, exposée via des API (\textit{full API}) et reposant sur une architecture orientée microservices.

Le projet étant issu d'une vision métier définie au sein de l'entreprise, son organisation repose sur un développement par modules. Cette approche permet de préciser progressivement les fonctionnalités et de les ajuster au fur et à mesure de leur conception.

Travailler dans un environnement bancaire impose un cadre exigeant. La sécurité, la conformité réglementaire (KYC), la traçabilité des opérations et l'ergonomie des interfaces sont autant de contraintes à respecter, auxquelles s'ajoute l'intégration avec le cœur du système bancaire existant.

Plusieurs risques ont par ailleurs été identifiés. Les besoins métiers peuvent être mal compris, l'intégration avec les systèmes existants s'avérer complexe, les données disponibles se révéler incomplètes ou les règles de conformité mal interprétées ; s'y ajoute le risque d'une adoption trop faible par les utilisateurs finaux.

Les résultats attendus concernent la mise en place d'un parcours client plus unifié, une meilleure centralisation de la connaissance client, un accès intégré aux services d'investissement et un appui au pilotage des activités métier. Le projet vise également à limiter les ruptures applicatives entre les différents systèmes.

Compte tenu de ces éléments, une approche entièrement figée dès le départ apparaît inadaptée. En effet, l'ensemble des fonctionnalités attendues ne peut pas toujours être défini avec précision au début du projet. Ainsi, les modules \textbf{RELATION CLIENT} et \textbf{SGI} sont traités de manière itérative, à travers des versions successives et une validation régulière des priorités, des parcours utilisateurs et des règles métiers.


% ─────────────────────────────────────────────────
\subsection{Méthodologie de travail}

Concrètement, le travail a démarré par un cadrage, pour poser le périmètre,
identifier les acteurs et s'entendre sur les objectifs et les contraintes. Une
fois ce cadre fixé, nous avons étudié les processus métier --- relation client,
KYC, gestion de portefeuille, passage d'ordres --- afin de comprendre ce que le
système devait réellement faire. Ces observations ont été traduites en modèles,
en nous appuyant sur \textbf{UML} pour représenter les fonctionnalités, les
interactions et les flux de données. La conception applicative a suivi, avec le
découpage en modules et microservices, puis la définition des API, des rôles et
des principales règles de gestion.

Plutôt que de tout figer au départ, nous avancions par étapes. Toutes les deux
semaines, le vendredi, nous présentions l'avancement du travail. De temps en
temps, des experts d'un module précis --- la SGI par exemple --- venaient
apporter leur regard métier, ce qui nous aidait à corriger et à améliorer ce qui
avait été développé. Les fonctionnalités étaient ainsi construites puis intégrées
au fur et à mesure au reste d'I-SIB.


\section*{Conclusion}
\addcontentsline{toc}{section}{Conclusion}

Ce chapitre a posé les bases du projet en présentant le contexte global dans
lequel il s'inscrit et les enjeux liés à la transformation des systèmes
d'information bancaires. La description de la structure d'accueil a permis de
mieux comprendre l'environnement professionnel, tandis que l'analyse du sujet a
mis en évidence les limites des solutions existantes et l'intérêt d'une approche
intégrée.

La problématique identifiée souligne l'importance de rapprocher les modules de relation client et de gestion d'investissement afin d'améliorer l'expérience utilisateur et le pilotage des activités. Les objectifs définis traduisent cette volonté à travers la mise en place d'une plateforme unifiée, intelligente et orientée données.

Enfin, la démarche méthodologique retenue, basée sur une approche itérative et
incrémentale, apparaît adaptée aux contraintes du projet et aux besoins
évolutifs du client. Elle favorise une meilleure maîtrise des risques, une
validation continue des fonctionnalités et une adéquation progressive entre la
solution développée et les attentes métiers.

Le chapitre suivant sera consacré à l'analyse détaillée des besoins et à la modélisation fonctionnelle du système proposé.
